\begin{table}[h!] \centering \caption{Summary Statistics China}\label{tab:summary_stats_CHN_unb} \begin{threeparttable} \begin{tabular}{l cccc} \toprule
& All & Rural & Urban & Difference \\
 &  &  &  & $t$-test \\
\midrule
Location &  &   54.2\% &   45.8\% &  \\  \addlinespace \addlinespace
Log Consumption & 10.46 & 10.26 & 10.70 & -0.44*** \\
 & (0.93) & (0.90) & (0.90) &  \\  \addlinespace
Log Income & 8.74 & 8.25 & 9.20 & -0.95*** \\
 & (1.93) & (2.08) & (1.64) &  \\  \addlinespace
Female & 0.51 & 0.51 & 0.52 & -0.01*** \\
 & (0.50) & (0.50) & (0.50) &  \\  \addlinespace
Age (years) & 46.66 & 46.70 & 46.60 & 0.10 \\
 & (16.57) & (16.51) & (16.64) &  \\  \addlinespace
Education (years) & 7.38 & 6.16 & 8.83 & -2.68*** \\
 & (4.92) & (4.68) & (4.80) &  \\  \addlinespace
Household Size & 4.26 & 4.56 & 3.91 & 0.65*** \\
 & (1.91) & (1.98) & (1.76) &  \\  \addlinespace


\cmidrule{2-5}
Observations & 109,535 & 59,354 & 50,181 &  \\  
Individuals & 34,746 &  &  &  \\
Non-switchers &   95.7\% &  &  &  \\  \addlinespace
\bottomrule \end{tabular} \begin{tablenotes}[flushleft] \footnotesize \item{Summary statistics for China for the unbalanced panel across all waves. Source: China survey. The table reports means and standard deviations (in parentheses) based on individual-year pairs. See section 3 for further details. All variables have the same number of observations, except for income, which is missing for some observations. Income has 49,191 observations.} \end{tablenotes} \end{threeparttable} \end{table}